\section{Αντικείμενο και κίνητρο της διπλωματικής}
Η λειτουργία μιας σύγχρονης διαδικτυακής εφαρμογής δεν εξαρτάται μόνο από
την ορθότητα του πηγαίου κώδικα. Εξαρτάται επίσης από την υπολογιστική
υποδομή στην οποία εκτελείται, τον τρόπο με τον οποίο αποθηκεύονται τα
δεδομένα, την προστασία της δικτυακής επικοινωνίας, τη δυνατότητα
παρακολούθησης της κατάστασης του συστήματος και την ορθή επανεκκίνηση των
services μετά από προγραμματισμένη διακοπή. Οι παράγοντες αυτοί αποκτούν ιδιαίτερη σημασία σε
εφαρμογές των οποίων η διακοπή ή η απώλεια δεδομένων επηρεάζει ουσιώδεις
επιχειρησιακές λειτουργίες.

Το υπολογιστικό νέφος επιλέχθηκε επειδή επιτρέπει την κατά απαίτηση διάθεση
υπολογιστικών, αποθηκευτικών και δικτυακών πόρων χωρίς την προμήθεια
ιδιόκτητου φυσικού εξοπλισμού. Η έννοια, τα μοντέλα υπηρεσιών και οι
βιβλιογραφικές της βάσεις αναπτύσσονται συστηματικά στο Κεφάλαιο~2. Στο
παρόν κεφάλαιο ενδιαφέρει κυρίως το πρακτικό κίνητρο: η δυνατότητα να
μετατραπεί μία τοπική εφαρμογή σε δημόσια, ασφαλή και μετρήσιμη υπηρεσία,
με ελεγχόμενο λειτουργικό κόστος.

Αντικείμενο της παρούσας διπλωματικής εργασίας είναι η μελέτη, υλοποίηση και
αξιολόγηση υποδομής cloud για τη φιλοξενία μιας ολοκληρωμένης διαδικτυακής
εφαρμογής. Ως μελέτη περίπτωσης χρησιμοποιείται το \emph{Happy Tails}, ένα
σύστημα διαχείρισης κτηνιατρικής κλινικής που υποστηρίζει ιδιοκτήτες,
κατοικίδια, ραντεβού, επισκέψεις, ιατρικό ιστορικό και σχετικά έγγραφα. Η
εφαρμογή συνδυάζει διεπαφή χρήστη σε React, backend σε Spring Boot και
σχεσιακή βάση δεδομένων MySQL \cite{react2026docs,spring2026boot,mysql2026manual}.
Επομένως αποτελεί κατάλληλο πραγματικό
φορτίο εργασίας για την εξέταση ζητημάτων ανάπτυξης, ασφάλειας, επίδοσης,
διαθεσιμότητας, συντηρησιμότητας και κόστους.

Η αφετηρία της εργασίας ήταν ένα πρακτικό πρόβλημα: η εφαρμογή λειτουργούσε
τοπικά, αλλά δεν υπήρχε μία πλήρως καταγεγραμμένη διαδρομή για τη δημόσια και
ασφαλή διάθεσή της. Η μετάβαση απαίτησε σταθερό domain name, κρυπτογραφημένη
επικοινωνία και διαδικασία deployment που να μπορεί να επαναληφθεί χωρίς να
εξαρτάται από χειροκίνητες, μη καταγεγραμμένες ενέργειες. Παράλληλα, ο
περιορισμένος προϋπολογισμός οδήγησε στην επιλογή μίας EC2 instance και
δημιούργησε το ουσιαστικό ερώτημα της μελέτης: ποια χαρακτηριστικά μπορούν να
καλυφθούν από αυτή την οικονομική λύση και ποια θα απαιτούσαν πρόσθετες
managed services. Η βιβλιογραφία για το cloud migration επιβεβαιώνει ότι η
διαδικασία δεν είναι απλή μεταφορά εκτελέσιμων αρχείων, αλλά απαιτεί
αρχιτεκτονική προσαρμογή και μεθοδική λήψη αποφάσεων
\cite{jamshidi2013migration}.

Ο όρος \emph{mission-critical} στον τίτλο εκφράζει το σύνολο των
σχεδιαστικών απαιτήσεων που εξετάζονται---συνέχεια λειτουργίας, προστασία
δεδομένων, ελεγχόμενη πρόσβαση, παρατηρησιμότητα και επαληθεύσιμη
επανεκκίνηση---και όχι ισχυρισμό ότι το Happy Tails αποτελεί πιστοποιημένο
σύστημα ιατρικής ασφάλειας. Η εφαρμογή λειτουργεί ως αντιπροσωπευτική
μελέτη περίπτωσης πάνω στην οποία εφαρμόζονται και αξιολογούνται αυτές οι
αρχές σε ρεαλιστικό, αλλά ελεγχόμενο, ακαδημαϊκό περιβάλλον.

\section{Ορισμός του προβλήματος και πεδίο εφαρμογής}
Η αρχική κατάσταση του συστήματος χαρακτηριζόταν από τη συνύπαρξη του
frontend, του backend και της βάσης δεδομένων σε περιβάλλον κατάλληλο για
ανάπτυξη και δοκιμές, χωρίς πλήρως τεκμηριωμένο μηχανισμό δημόσιας διάθεσης
και λειτουργικής παρακολούθησης. Για να μετατραπεί το σύστημα σε αξιόπιστη
διαδικτυακή υπηρεσία απαιτούνται απαντήσεις σε πρακτικά ερωτήματα: ποιοι
υπολογιστικοί πόροι είναι επαρκείς, ποια ports πρέπει να είναι δημόσια,
πώς διατηρείται σταθερή η διεύθυνση της υπηρεσίας, πώς εκδίδεται και
ανανεώνεται ένα TLS certificate, πώς επανεκκινούν τα επιμέρους components
και πώς αποδεικνύεται ότι η εφαρμογή λειτουργεί μετά από αλλαγή ή αστοχία.

Για τις ανάγκες της εργασίας, η αξιοπιστία μεταφράζεται σε συγκεκριμένες
τεχνικές απαιτήσεις: ορθή εκκίνηση των υπηρεσιών, προστασία της επικοινωνίας,
περιορισμό της διαχειριστικής πρόσβασης, διατήρηση των δεδομένων, καταγραφή
της κατάστασης και δημιουργία αντιγράφων ασφαλείας. Η θεωρητική διάκριση
μεταξύ reliability, availability, integrity και maintainability παρουσιάζεται
στο Κεφάλαιο~2, ώστε εδώ να παραμένει καθαρός ο ορισμός του πρακτικού
προβλήματος.

Το πεδίο εφαρμογής περιλαμβάνει την ανάλυση του Happy Tails, τη συσκευασία
των components του σε containers, την ανάπτυξή του στο AWS, τη ρύθμιση
δικτύου και DNS, την ασφαλή δημοσίευση μέσω HTTPS και την πειραματική
αξιολόγηση της τελικής εγκατάστασης. Η ενεργή υλοποίηση βασίζεται σε EC2,
EBS, Elastic IP και Route~53, σε συνδυασμό με Docker Compose, Nginx, Spring
Boot και MySQL. Υπηρεσίες όπως Amazon RDS/Aurora, Elastic Load Balancing,
Auto Scaling και S3 εξετάζονται ως εναλλακτικές ή μελλοντικές
αρχιτεκτονικές. Το CloudFront αξιολογείται χωριστά ως περιορισμένο proof of
concept, χωρίς μεταβολή του production DNS ή ένταξή του στη βασική διαδρομή
της εφαρμογής.
Σε όλο τον τόμο χρησιμοποιούνται τρεις σαφείς χαρακτηρισμοί:
\emph{υλοποιημένο}, όταν υπάρχει ενεργή διαμόρφωση και τεχνικό τεκμήριο,
\emph{σε αναμονή επαλήθευσης}, όταν έχει ολοκληρωθεί ο κώδικας ή η προετοιμασία
αλλά όχι η τελική δοκιμή, και \emph{προτεινόμενο}, όταν πρόκειται για
αρχιτεκτονική επιλογή που αξιολογείται μόνο θεωρητικά.

Η ποιότητα της εφαρμογής και της υποδομής αντιμετωπίζεται ως μετρήσιμο
αντικείμενο. Τα κριτήρια οργανώνονται με βάση το ISO/IEC~25010:2023 και
εξειδικεύονται ως λειτουργικές και μη λειτουργικές απαιτήσεις στο
Κεφάλαιο~3. Δεν επιχειρείται πιστοποίηση κατά το πρότυπο· χρησιμοποιείται
μόνο ως πλαίσιο για τη σύνδεση απαιτήσεων, δοκιμών και κριτηρίων αποδοχής
\cite{iso25010-2023}.

\section{Στόχοι και ερευνητικά ερωτήματα}
Κύριος στόχος είναι η σχεδίαση και τεκμηριωμένη αξιολόγηση μιας λειτουργικής
cloud υποδομής για το Happy Tails. Ο γενικός αυτός στόχος αναλύεται στους
ακόλουθους επιμέρους στόχους:

\begin{enumerate}[label=\arabic*.]
  \item καταγραφή των λειτουργικών και μη λειτουργικών απαιτήσεων της
        εφαρμογής και χαρτογράφηση των components και των εξαρτήσεών της,
  \item επαναλήψιμη ανάπτυξη του συστήματος σε AWS EC2 με διατήρηση δεδομένων
        και σαφή διαχωρισμό εφαρμογής, database και reverse proxy,
  \item ασφαλής δημόσια πρόσβαση μέσω περιορισμένων security-group rules,
        σταθερής Elastic IP, Route~53 και HTTPS,
  \item επαλήθευση της λειτουργίας με τεχνικούς ελέγχους και
        αντιπροσωπευτικές end-to-end ροές χρήσης,
  \item μέτρηση χρόνου απόκρισης, χρήσης πόρων και χρόνου επανεκκίνησης σε
        ελεγχόμενα σενάρια,
  \item αξιολόγηση κόστους, περιορισμών και δυνατοτήτων εξέλιξης προς
        αρχιτεκτονική μεγαλύτερης διαθεσιμότητας.
\end{enumerate}

Με βάση τους στόχους διατυπώνονται τα ακόλουθα ερευνητικά ερωτήματα:

\begin{description}[style=nextline,leftmargin=1.6cm,labelwidth=1.2cm]
  \item[ΕΕ1:] Σε ποιο βαθμό μπορεί μία οικονομικά περιορισμένη, containerized
  εγκατάσταση σε EC2 να υποστηρίξει με ασφαλή και επαναλήψιμο τρόπο τις
  βασικές λειτουργίες μιας εφαρμογής διαχείρισης κτηνιατρικής κλινικής;

  \item[ΕΕ2:] Ποια είναι η παρατηρούμενη συμπεριφορά της εγκατάστασης ως προς
  τους χρόνους απόκρισης, τη χρήση CPU και μνήμης, την ορθή εκκίνηση και τον
  χρόνο επανεκκίνησης των services μετά από ελεγχόμενο reboot του instance;

  \item[ΕΕ3:] Ποιοι περιορισμοί και single points of failure παραμένουν στην
  υλοποιημένη αρχιτεκτονική και ποιες AWS managed services θα μπορούσαν να
  τους αντιμετωπίσουν;

  \item[ΕΕ4:] Ποιο είναι το κόστος και ποιοι οι βασικοί συμβιβασμοί μεταξύ
  της υλοποιημένης λύσης και μιας προτεινόμενης αρχιτεκτονικής υψηλότερης
  διαθεσιμότητας;
\end{description}

Για κάθε ερευνητικό ερώτημα απαιτείται παρατηρήσιμο αποτέλεσμα και όχι μόνο
περιγραφική αναφορά μιας υπηρεσίας. Τα τεκμήρια περιλαμβάνουν πίνακες
παραμετροποίησης, status checks, health checks, logs, HTTP status codes,
ολοκληρωμένες ροές της εφαρμογής και μετρήσεις πόρων. Ένα στιγμιότυπο της
αρχικής σελίδας αποδεικνύει ότι η διεπαφή είναι προσβάσιμη, όχι όμως ότι η
υποδομή λειτουργεί αξιόπιστα ως σύνολο.

\section{Μεθοδολογική προσέγγιση}
Η εργασία ακολουθεί μεθοδολογία μελέτης περίπτωσης και πειραματικής
αξιολόγησης. Πρώτο αντικείμενο μελέτης είναι ο ίδιος ο πηγαίος κώδικας του
Happy Tails: οι ρόλοι χρηστών, οι κύριες ροές, τα δεδομένα και οι εξαρτήσεις
του. Από αυτή την καταγραφή προκύπτουν οι απαιτήσεις της υποδομής και τα
κριτήρια με τα οποία ελέγχεται κάθε αλλαγή.

Η cloud αρχιτεκτονική αναπτύσσεται σταδιακά και κάθε ουσιώδης επιλογή
συνοδεύεται από αιτιολόγηση και έλεγχο. Αυτό ισχύει για το instance type, το
storage, το security group, το DNS, τον reverse proxy και το TLS certificate.
Η τεκμηρίωση συνδυάζει redacted screenshots, πίνακες παραμετροποίησης, logs
και επαναλήψιμες εντολές. Έτσι, η περιγραφή μιας υπηρεσίας AWS παραμένει
διακριτή από το τεκμήριο ότι η υπηρεσία χρησιμοποιήθηκε στη συγκεκριμένη
εγκατάσταση.

Η αξιολόγηση θα περιλάβει λειτουργικούς ελέγχους του Happy Tails, επαλήθευση
της persistence των δεδομένων, ελεγχόμενη επανεκκίνηση της EC2 instance,
καταγραφή του recovery time και βασικές δοκιμές φορτίου. Οι μετρήσεις θα
εκτελούνται με καταγεγραμμένες αρχικές συνθήκες, συγκεκριμένο πλήθος
αιτημάτων και σαφή εργαλεία, ώστε τα αποτελέσματα να είναι αναπαραγώγιμα.
Οι μετρικές θα συσχετιστούν με το instance type και το πραγματικό κόστος της
περιόδου λειτουργίας.

Οι αρχιτεκτονικές αποφάσεις ελέγχονται επιπλέον με τα κριτήρια του AWS
Well-Architected Framework, τα οποία παρουσιάζονται στο Κεφάλαιο~2
\cite{aws2024wellarchitected}. Το πλαίσιο λειτουργεί ως δομημένος κατάλογος
ερωτήσεων και όχι ως μηχανισμός πιστοποίησης. Έτσι, η αξιολόγηση δεν
περιορίζεται στο αν η εφαρμογή εκτελείται, αλλά εξετάζει και τις συνθήκες
λειτουργίας, ασφάλειας και συντήρησής της.

\section{Συνεισφορά}
Η επιδιωκόμενη συνεισφορά της εργασίας είναι η σύνδεση μιας πραγματικής
full-stack εφαρμογής με μία τεκμηριωμένη διαδικασία cloud deployment και
αξιολόγησης. Στο παρόν στάδιο έχουν ολοκληρωθεί η βασική εγκατάσταση και οι
αρχικοί λειτουργικοί έλεγχοι, ενώ οι μετρήσεις επίδοσης, επανεκκίνησης και
κόστους θα οριστικοποιήσουν το πειραματικό μέρος. Η εργασία περιλαμβάνει:

\begin{itemize}
  \item ανάλυση της αρχιτεκτονικής και των βασικών ροών του Happy Tails,
  \item λειτουργική containerized εγκατάσταση Spring Boot, MySQL και Nginx
        σε AWS EC2,
  \item τεκμηρίωση της δημιουργίας του EC2 instance, της δικτυακής προστασίας,
        του DNS και του HTTPS,
  \item σύνολο επαληθεύσεων που συνδέει την κατάσταση της υποδομής με
        πραγματικές λειτουργίες της εφαρμογής,
  \item σχέδιο και, μετά την ολοκλήρωση των μετρήσεων, πειραματική αποτίμηση
        επίδοσης, επανεκκίνησης και κόστους,
  \item κριτική σύγκριση της υλοποιημένης οικονομικής λύσης με μία
        προτεινόμενη αρχιτεκτονική που αξιοποιεί managed services.
\end{itemize}

Το ενδιαφέρον της μελέτης δεν βρίσκεται σε μία γενική παρουσίαση των AWS
services. Βρίσκεται στις συγκεκριμένες επιλογές, στους περιορισμούς που
παρατηρήθηκαν και στην αντιστοίχιση κάθε τεχνικού ισχυρισμού με τεκμήριο ή
βιβλιογραφική πηγή.

\section{Οργάνωση του τόμου}
Ο υπόλοιπος τόμος οργανώνεται σε έξι κεφάλαια. Το Κεφάλαιο~2 παρουσιάζει τη
βιβλιογραφική ανασκόπηση και το θεωρητικό υπόβαθρο του cloud computing, της
αξιοπιστίας, των AWS services και των cloud-native εφαρμογών. Το
Κεφάλαιο~3 αναλύει τις απαιτήσεις, την αρχιτεκτονική και τις βασικές ροές
του Happy Tails. Το Κεφάλαιο~4 διατυπώνει τους σχεδιαστικούς στόχους και
παρουσιάζει τον σχεδιασμό της cloud υποδομής, μαζί με τις εναλλακτικές
αρχιτεκτονικές που εξετάστηκαν. Το Κεφάλαιο~5 περιγράφει την πραγματική
υλοποίηση και τη διαδικασία deployment στο AWS. Το Κεφάλαιο~6 παρουσιάζει
τη μεθοδολογία και τα αποτελέσματα της πειραματικής αξιολόγησης. Τέλος, το
Κεφάλαιο~7 συνοψίζει τα συμπεράσματα, τους περιορισμούς και τις προτεινόμενες
μελλοντικές επεκτάσεις.
